\chapter{Sylow Theorems and Finite Groups}

\begin{center}\itshape
Is there a simple group of order $30$?\\
Lagrange says what subgroup orders are allowed; Sylow says which ones must occur, and how many times.\\
Counting them answers the question in three lines.
\end{center}

Ludwig Sylow taught mathematics at a secondary school in Halden, Norway, for forty years. In 1872 he published, in a ten-page paper, the
theorems that bear his name, and they became the first tools of a project that took a century: the classification of all finite simple groups,
the ``atoms'' of finite group theory. The project ended with a list, and the list has a famous last entry, the \emph{Monster}, with about
$8\cdot10^{53}$ elements. Everything there starts with the counting in this chapter.

Lagrange's theorem forbids: a group of order $12$ has no subgroup of order $5$. It promises nothing: $A_4$ has no subgroup of order $6$ (Proposition~\ref{prop:a4}). Cauchy's theorem (Theorem~\ref{thm:cauchy}) promises subgroups of prime order. The three theorems of Sylow (1872) go much further: if $p^a$ is the largest power of $p$ dividing $|G|$, there are subgroups of order $p^a$, they are all conjugate, and their number is $\equiv1\pmod p$. With these, one can classify the groups of small order, prove that many orders admit no simple group, and set up the notion of a \emph{solvable} group, which Galois theory needs (Chapter~10).

Throughout, $G$ is a finite group, $p$ a prime, and $|G|=p^am$ with $p\nmid m$.

\section{The three theorems}

\begin{definition}
A \emph{$p$-subgroup} of $G$ is a subgroup whose order is a power of $p$. A \emph{Sylow $p$-subgroup} is a subgroup of order $p^a$. Their set is $\operatorname{Syl}_p(G)$ and their number $n_p$. The \emph{normaliser} of a subgroup $H$ is $N_G(H)=\{g\in G:gHg^{-1}=H\}$.
\end{definition}

$N_G(H)$ is the stabiliser of $H$ for the action of $G$ on its subgroups by conjugation, so it is a subgroup, and $H\trianglelefteq N_G(H)$; it is the largest subgroup of $G$ in which $H$ is normal. By the orbit--stabiliser theorem, $H$ has $[G:N_G(H)]$ conjugates.

\begin{lemma}\label{lem:binom}
$\binom{p^am}{p^a}\equiv m\pmod p$. In particular it is not divisible by $p$.
\end{lemma}
\begin{proof}
In $(\Z/p\Z)[X]$, $(1+X)^p=1+X^p$, because $p$ divides $\binom pk$ for $0<k<p$ (Tome~1). Raising to the power $p$ repeatedly, $(1+X)^{p^a}=1+X^{p^a}$, so
\[
(1+X)^{p^am}=\bigl(1+X^{p^a}\bigr)^m=\sum_{j=0}^m\binom mjX^{jp^a}.
\]
Compare the coefficients of $X^{p^a}$: on the left $\binom{p^am}{p^a}$ reduced mod $p$, on the right $\binom m1=m$.
\end{proof}

\begin{theorem}[Sylow I]\label{thm:sylow1}
$G$ has a subgroup of order $p^a$.
\end{theorem}
\begin{proof}[Proof (Wielandt, 1959)]
Let $\Omega$ be the set of subsets of $G$ with $p^a$ elements. $G$ acts on $\Omega$ by left multiplication, $g\cdot A=gA$. By Lemma~\ref{lem:binom}, $|\Omega|$ is not divisible by $p$, and $\Omega$ is the disjoint union of the orbits, so some orbit $O$ has size prime to $p$. Pick $A\in O$ and let $H=\{g:gA=A\}$ be its stabiliser.
\begin{itemize}
\item $|O|=|G|/|H|$ (Theorem~\ref{thm:orbstab}) is prime to $p$, so $p^a$ divides $|H|$.
\item Fix $x\in A$. For $h\in H$, $hx\in hA=A$, and $h\mapsto hx$ is injective. So $|H|\le|A|=p^a$.
\end{itemize}
Hence $|H|=p^a$.
\end{proof}

\begin{think}
Why would all Sylow $p$-subgroups be conjugate? Try $\mathfrak S_3$ and $p=2$: list its subgroups of order $2$ and conjugate one of them. Then guess $n_2$ and $n_3$ for $\mathfrak S_4$.
\end{think}

\begin{theorem}[Sylow II and III]\label{thm:sylow23}
Let $P$ be a Sylow $p$-subgroup of $G$.
\begin{enumerate}[label=(\alph*)]
\item Every $p$-subgroup of $G$ is contained in a conjugate of $P$.
\item The Sylow $p$-subgroups are exactly the conjugates of $P$.
\item $n_p=[G:N_G(P)]$, so $n_p$ divides $m$; and $n_p\equiv1\pmod p$.
\end{enumerate}
\end{theorem}
\begin{proof}
(a) Let $Q$ be a $p$-subgroup. $Q$ acts on the set $G/P$ of left cosets by left multiplication. There are $m$ cosets, prime to $p$, and each orbit has size $[Q:Q_x]$, a power of $p$; so some orbit has size $1$, i.e.\ some coset $gP$ satisfies $qgP=gP$ for all $q\in Q$. Then $g^{-1}qg\in P$ for all $q\in Q$, i.e.\ $Q\subseteq gPg^{-1}$.

(b) Conjugates of $P$ have order $p^a$. Conversely, if $Q$ is a Sylow $p$-subgroup, by (a) $Q\subseteq gPg^{-1}$, and both have order $p^a$, so $Q=gPg^{-1}$.

(c) By (b), $G$ acts transitively on $\operatorname{Syl}_p(G)$ by conjugation, and the stabiliser of $P$ is $N_G(P)$; so $n_p=[G:N_G(P)]$. Since $P\le N_G(P)$, $n_p$ divides $[G:P]=m$.

Now let $P$ itself act on $\operatorname{Syl}_p(G)$ by conjugation. Orbits have size a power of $p$. A fixed point is a Sylow subgroup $Q$ with $xQx^{-1}=Q$ for all $x\in P$, i.e.\ $P\subseteq N_G(Q)$. Then $P$ and $Q$ are both Sylow $p$-subgroups of the group $N_G(Q)$ (their order $p^a$ is the largest possible already in $G$), so by (b) applied to $N_G(Q)$ they are conjugate in $N_G(Q)$. But $Q$ is normal in $N_G(Q)$, so its only conjugate there is itself: $P=Q$. So $P$ is the only fixed point, every other orbit has size divisible by $p$, and $n_p\equiv1\pmod p$.
\end{proof}

\begin{corollary}\label{cor:sylownormal}
A Sylow $p$-subgroup $P$ is normal iff $n_p=1$. In that case $P$ contains every $p$-subgroup of $G$, and every element of order a power of $p$.
\end{corollary}
\begin{proof}
$P$ is normal iff it equals all its conjugates, iff $n_p=1$ by (b). Then (a) puts every $p$-subgroup, in particular every $\langle g\rangle$ with $g$ of order $p^k$, inside $P$.
\end{proof}

\begin{corollary}\label{cor:pk}
If $p^k$ divides $|G|$, then $G$ has a subgroup of order $p^k$.
\end{corollary}
\begin{proof}
By Theorem~\ref{thm:sylow1} it suffices to show that a group $P$ of order $p^a$ has a subgroup of each order $p^k$, $0\le k\le a$. Induction on $a$; the case $k=0$ is clear. For $k\ge1$, the centre $Z(P)$ is non-trivial (Corollary~\ref{cor:pgroup}), so by Cauchy's theorem it contains an element of order $p$, which generates a subgroup $C$ of order $p$, normal because central. $P/C$ has order $p^{a-1}$, so by induction it has a subgroup of order $p^{k-1}$, which is $K/C$ for some $C\le K\le P$ (Theorem~\ref{thm:iso23}(a)); then $|K|=p^k$.
\end{proof}

\section{Using Sylow}

The method: factor $|G|$, list for each $p$ the divisors of $m$ that are $\equiv1\pmod p$, and count elements.

\begin{proposition}\label{prop:thirty}
No group of order $30$ is simple.
\end{proposition}
\begin{proof}
$30=2\cdot3\cdot5$. $n_5$ divides $6$ and is $\equiv1\pmod5$: $n_5\in\{1,6\}$. $n_3$ divides $10$ and is $\equiv1\pmod3$: $n_3\in\{1,10\}$. If neither is $1$, there are $6$ subgroups of order $5$, which pairwise meet in $\{e\}$ (their intersection is a subgroup of a group of prime order), so they contain $6\cdot4=24$ elements of order $5$; likewise $10\cdot2=20$ elements of order $3$. That is $44>30$ elements. So $n_5=1$ or $n_3=1$, and that Sylow subgroup is a proper non-trivial normal subgroup (Corollary~\ref{cor:sylownormal}).
\end{proof}

\begin{theorem}[groups of order $pq$]\label{thm:pq}
Let $p<q$ be primes and $|G|=pq$.
\begin{enumerate}[label=(\alph*)]
\item If $p\nmid q-1$, then $G\cong\Z/pq\Z$.
\item If $p\mid q-1$, then $G$ is either cyclic or isomorphic to the non-abelian semidirect product $\Z/q\Z\rtimes_u\Z/p\Z$, where $b\in\Z/p\Z$ acts on $\Z/q\Z$ as multiplication by $u^b$ and $u$ is any element of order $p$ in $(\Z/q\Z)^*$; this group does not depend on the choice of $u$.
\end{enumerate}
In particular every group of order $2q$ ($q$ an odd prime) is $\Z/2q\Z$ or the dihedral group $D_q$.
\end{theorem}
\begin{proof}
$n_q$ divides $p$ and is $\equiv1\pmod q$; since $p<q$, $n_q=1$: the Sylow $q$-subgroup $Q=\langle a\rangle$ is normal. Let $P=\langle b\rangle$ be a Sylow $p$-subgroup. $P\cap Q=\{e\}$ (orders are coprime) and $|PQ|=|P||Q|/|P\cap Q|=pq$ (the products $xy$, $x\in P$, $y\in Q$, are distinct), so by Proposition~\ref{prop:semidirect}(b), $G\cong Q\rtimes P$ with $b$ acting by $bab^{-1}=a^u$ for some $u\in(\Z/q\Z)^*$ with $u^p=1$ (since $b^p=e$).

(a) $n_p$ divides $q$ and is $\equiv1\pmod p$; as $q\not\equiv1\pmod p$, $n_p=1$. Both $P$ and $Q$ are normal, so $G\cong P\times Q\cong\Z/p\Z\times\Z/q\Z\cong\Z/pq\Z$ (Proposition~\ref{prop:internal} and Theorem~\ref{thm:crt}).

(b) If $u=1$, $G$ is abelian and cyclic as in (a). Otherwise $u$ has order $p$. The group $(\Z/q\Z)^*$ is cyclic of order $q-1$ (primitive roots, Tome~1), so it has exactly one subgroup of order $p$ (Theorem~\ref{thm:cyclic}), and the elements of order $p$ are $u,u^2,\dots,u^{p-1}$. Replacing the generator $b$ of $P$ by $b^k$ ($p\nmid k$) replaces $u$ by $u^k$. So whatever the element of order $p$, after choosing the right generator of $P$ the multiplication rule of $G$ in terms of $a$, $b$ is the same: all these groups are isomorphic. Such a group exists (Proposition~\ref{prop:semidirect}(a)) and is not abelian, since $bab^{-1}=a^u\ne a$.

For $p=2$: $2$ divides $q-1$, and the only element of order $2$ in $(\Z/q\Z)^*$ is $-1$: $bab^{-1}=a^{-1}$, which is the rule of $D_q$ (rotation $a$, reflection $b$).
\end{proof}

So a group of order $15$, $33$, $35$, $51$ or $65$ is cyclic, while order $21$ has exactly two groups, $\Z/21\Z$ and $\Z/7\Z\rtimes\Z/3\Z$ (the group of the notebook of Chapter~2). Together with Exercise~\ref{ex:p2}, Exercise~\ref{ex:order8} and Exercise~\ref{ex:order12} below, this gives the complete list of groups of order at most $15$, in Figure~\ref{fig:small}.

\begin{figure}[ht]
\centering
\fig{ch05-small}
\caption{All groups of order $1$ to $15$ up to isomorphism, abelian ones in blue and non-abelian ones in orange ($\Z_n$ stands for $\Z/n\Z$). ``Prime'' is Corollary~\ref{cor:lagrange}, ``$p^2$'' is Exercise~\ref{ex:p2}, ``$2p$'' and ``$pq$'' are Theorem~\ref{thm:pq}, order $8$ is Exercise~\ref{ex:order8} and order $12$ is Exercise~\ref{ex:order12}. Order $16$ has $14$ groups, and order $2^{10}=1024$ has $49\,487\,365\,422$: the $p$-groups are by far the most numerous.}
\label{fig:small}
\end{figure}

\begin{think}
For $|G|=12$: list the possible values of $n_2$ and $n_3$. Which combinations can occur together? Which group has $n_3=4$?
\end{think}

\section{Solvable groups}

\begin{definition}
$G$ is \emph{solvable} if there is a chain of subgroups
\[
\{e\}=G_0\trianglelefteq G_1\trianglelefteq\dots\trianglelefteq G_r=G,
\]
each normal in the next (not necessarily in $G$), with every quotient $G_i/G_{i-1}$ abelian.
\end{definition}

The name comes from Galois theory: a polynomial equation is solvable by radicals iff its Galois group is solvable (Chapter~10). Abelian groups are solvable ($r=1$), and so is $\mathfrak S_4$:
\[
\{e\}\trianglelefteq V\trianglelefteq A_4\trianglelefteq\mathfrak S_4,
\]
with $V$ the Klein group of Exercise~\ref{ex:v4} and quotients $V$, $A_4/V\cong\Z/3\Z$, $\mathfrak S_4/A_4\cong\Z/2\Z$, all abelian.

\begin{proposition}\label{prop:solv}
\begin{enumerate}[label=(\alph*)]
\item Subgroups and quotients of a solvable group are solvable.
\item If $N\trianglelefteq G$ and both $N$ and $G/N$ are solvable, then $G$ is solvable.
\item Every $p$-group is solvable. Every group of order $pq$ is solvable.
\item A non-abelian simple group is not solvable. In particular $A_n$ and $\mathfrak S_n$ are not solvable for $n\ge5$.
\end{enumerate}
\end{proposition}
\begin{proof}
(a) Let $(G_i)$ be a chain as in the definition and $H\le G$. Put $H_i=H\cap G_i$. If $x\in H_i$, conjugation by $x$ preserves $H$ and $G_{i-1}$, so $H_{i-1}\trianglelefteq H_i$. The map $H_i\to G_i/G_{i-1}$, $x\mapsto xG_{i-1}$, is a homomorphism with kernel $H_i\cap G_{i-1}=H_{i-1}$, so $H_i/H_{i-1}$ is isomorphic to a subgroup of the abelian group $G_i/G_{i-1}$ (Theorem~\ref{thm:iso1}).

For $N\trianglelefteq G$ and $\pi:G\to G/N$, put $K_i=\pi(G_i)$. Since $G_{i-1}\trianglelefteq G_i$, $\pi(G_{i-1})\trianglelefteq\pi(G_i)$. The composite $G_i\to K_i\to K_i/K_{i-1}$ is onto, and $G_{i-1}$ lies in its kernel, so it induces a surjection $G_i/G_{i-1}\to K_i/K_{i-1}$; an image of an abelian group is abelian.

(b) Let $\{e\}=N_0\trianglelefteq\dots\trianglelefteq N_s=N$ and $\{e\}=K_0\trianglelefteq\dots\trianglelefteq K_r=G/N$ be chains with abelian quotients, and let $L_i=\pi^{-1}(K_i)$. Then $L_0=N$, $L_r=G$, $L_{i-1}\trianglelefteq L_i$ (the preimage of a normal subgroup under the surjection $L_i\to K_i$ is normal), and by the third isomorphism theorem $L_i/L_{i-1}\cong(L_i/N)/(L_{i-1}/N)=K_i/K_{i-1}$, abelian. The chain $N_0,\dots,N_s=L_0,L_1,\dots,L_r$ works for $G$.

(c) By induction on $|P|$: if $P\ne\{e\}$ is a $p$-group, its centre $Z$ is non-trivial and abelian, and $P/Z$ is a smaller $p$-group, solvable by induction; apply (b). For order $pq$, the normal Sylow $q$-subgroup $Q$ of the proof of Theorem~\ref{thm:pq} is cyclic, and $G/Q$ has order $p$: apply (b).

(d) Let $G$ be simple, non-abelian, with a chain as in the definition; delete repetitions, so that $G_{r-1}\ne G$. As $G_{r-1}\trianglelefteq G$ and $G$ is simple, $G_{r-1}=\{e\}$, so $G\cong G/G_{r-1}$ is abelian: a contradiction. So $A_n$ ($n\ge5$, Theorem~\ref{thm:an}) is not solvable, and neither is $\mathfrak S_n$, by (a).
\end{proof}

\begin{remark}
Two deep theorems extend (c). Burnside (1904): a group of order $p^\alpha q^\beta$ is solvable; the classical proof uses characters (Chapter~11). Feit and Thompson (1963): every group of odd order is solvable; their proof takes $255$ pages. One can also check, with the methods of this chapter but at much greater length, that every group of order less than $60$ is solvable, and that $A_5$ is the only simple group of order $60$; we will not need these facts.
\end{remark}

\begin{beauty}
The quintic in one picture: the roots of $x^5-4x+2$ are permuted by a group isomorphic to $\mathfrak S_5$ (Chapter~10), and $\mathfrak S_5$ is not solvable because $A_5$ is simple. The cube's rotations ($\mathfrak S_4$) are solvable, and this is why the quartic has a formula. Sylow's theorems also run modern computation: the discrete logarithm in a group of order $n$ is only as hard as in its largest Sylow subgroup (Pohlig--Hellman, in the notebook), which is why cryptography uses groups of prime order.
\end{beauty}

\section{Exercises}

\begin{exercise}[$\star$]\label{ex:sylowS4}
Find $n_2$ and $n_3$ for $\mathfrak S_4$, and describe its Sylow subgroups.
\end{exercise}

\begin{exercise}[$\star$]\label{ex:order45}
Show that every group of order $45$ is abelian. Which abelian groups of order $45$ are there?
\end{exercise}

\begin{exercise}[$\star\star$]\label{ex:nosimple}
Show that there is no simple group of order (a) $56$; (b) $200$; (c) $36$. For (c), let $G$ act on its Sylow $3$-subgroups.
\end{exercise}

\begin{exercise}[$\star\star$]\label{ex:S5sub}
Show that $\mathfrak S_5$ has no subgroup of order $15$, $30$ or $40$, although these numbers divide $120$.
\end{exercise}

\begin{exercise}[$\star\star$]\label{ex:normalizer}
Let $P$ be a Sylow $p$-subgroup of $G$ and $M=N_G(P)$. Show that $N_G(M)=M$.
\end{exercise}

\begin{exercise}[$\star\star$]\label{ex:derived}
With the derived subgroup of Exercise~\ref{ex:commutator}, define $G^{(0)}=G$ and $G^{(i+1)}=\bigl(G^{(i)}\bigr)'$.
\begin{enumerate}[label=(\alph*)]
\item Show that $G$ is solvable iff $G^{(k)}=\{e\}$ for some $k$.
\item Compute the derived series of $\mathfrak S_4$ and of $\mathfrak S_5$.
\end{enumerate}
\end{exercise}

\begin{exercise}[$\star\star$]\label{ex:sylowGL2}
Let $p$ be prime and $G=GL_2(\Z/p\Z)$, of order $(p^2-1)(p^2-p)$ (Exercise~\ref{ex:gl2}). Show that the matrices $\begin{psmallmatrix}1&x\\0&1\end{psmallmatrix}$ form a Sylow $p$-subgroup $U$, that $N_G(U)$ is the group $B$ of invertible upper triangular matrices, and that $n_p=p+1$. Interpret the $p+1$ Sylow subgroups geometrically.
\end{exercise}

\begin{exercise}[$\star\star\star$]\label{ex:order12}
Show that there are exactly five groups of order $12$ up to isomorphism: $\Z/12\Z$, $\Z/2\Z\times\Z/6\Z$, $A_4$, $D_6$ and $\Z/3\Z\rtimes\Z/4\Z$ (where a generator of $\Z/4\Z$ acts on $\Z/3\Z$ by $x\mapsto-x$).
\end{exercise}

\section{Hints}

\begin{hint}{ex:sylowS4}
$n_3$: count the $3$-cycles. $n_2$: use the action of $\mathfrak S_4$ on the three splittings into pairs (Exercise~\ref{ex:v4}).
\end{hint}
\begin{hint}{ex:order45}
Both Sylow subgroups are normal. Then use Exercise~\ref{ex:p2}.
\end{hint}
\begin{hint}{ex:nosimple}
(a) If $n_7=8$, count the elements of order $7$; what is left? (b) $n_5$ divides $8$. (c) $n_3\in\{1,4\}$; a homomorphism $G\to\mathfrak S_4$ cannot be injective.
\end{hint}
\begin{hint}{ex:S5sub}
For a subgroup $H$ of index $k$, let $\mathfrak S_5$ act on the $k$ cosets of $H$; the kernel is a normal subgroup contained in $H$. Use Exercise~\ref{ex:normalSn}. For $15$, use Theorem~\ref{thm:pq} and Exercise~\ref{ex:landau}.
\end{hint}
\begin{hint}{ex:normalizer}
If $g$ normalises $M$, then $gPg^{-1}$ is a Sylow $p$-subgroup of $M$. How many does $M$ have?
\end{hint}
\begin{hint}{ex:derived}
(a) If $G_r/G_{r-1}$ is abelian, $G'\subseteq G_{r-1}$; iterate. (b) The commutator of $(a\,b)$ and $(a\,c)$ is a $3$-cycle; in $A_4$, the quotient by $V$ is abelian.
\end{hint}
\begin{hint}{ex:sylowGL2}
The vectors fixed by every element of $U$ form the line spanned by $e_1$. If $gUg^{-1}=U$, where does $g$ send this line?
\end{hint}
\begin{hint}{ex:order12}
If $n_3=4$, let $G$ act on its four Sylow $3$-subgroups. If $n_3=1$, $G$ is a semidirect product $\Z/3\Z\rtimes Q$ with $|Q|=4$, and $\operatorname{Aut}(\Z/3\Z)$ has order $2$.
\end{hint}

\section{Solutions}

\begin{solution}{ex:sylowS4}
\idea Sylow subgroups of order $3$ are generated by elements of order $3$; for order $8$, find one and count its conjugates.
\writeup $|\mathfrak S_4|=24=2^3\cdot3$. A Sylow $3$-subgroup is cyclic of order $3$, generated by a $3$-cycle and containing two of them; there are $8$ three-cycles, so $n_3=4$ (and indeed $4\equiv1\pmod3$, $4\mid8$). For $p=2$, $\mathfrak S_4$ acts transitively on the three splittings $P_1,P_2,P_3$ into pairs (Exercise~\ref{ex:v4}), so the stabiliser $S_i$ of $P_i$ has order $24/3=8$: it is a Sylow $2$-subgroup. These three subgroups are distinct, since $(1\,2)$ fixes $P_1=\{12|34\}$ but not $P_2$. So $n_2\ge3$, and $n_2$ divides $3$: $n_2=3$. Each $S_i$ is a dihedral group $D_4$: number the corners of a square so that the two pairs of $P_i$ are the diagonals; its $8$ symmetries preserve the pair of diagonals.
\end{solution}

\begin{solution}{ex:order45}
\idea When all Sylow subgroups are normal, the group is their direct product.
\writeup $45=3^2\cdot5$. $n_5$ divides $9$ and is $\equiv1\pmod5$: $n_5=1$. $n_3$ divides $5$ and is $\equiv1\pmod3$: $n_3=1$. So the Sylow subgroups $P_3$ (order $9$) and $P_5$ (order $5$) are normal, meet trivially, and $|P_3P_5|=45$; by Proposition~\ref{prop:internal}, $G\cong P_3\times P_5$. $P_3$ is abelian by Exercise~\ref{ex:p2}, $P_5$ is cyclic, so $G$ is abelian, isomorphic to $\Z/9\Z\times\Z/5\Z\cong\Z/45\Z$ or $\Z/3\Z\times\Z/3\Z\times\Z/5\Z\cong\Z/3\Z\times\Z/15\Z$ (Theorem~\ref{thm:crt}). These two are not isomorphic: only the first has an element of order $9$.
\end{solution}

\begin{solution}{ex:nosimple}
\idea Count elements of prime order, or let the group act on a small set of Sylow subgroups to get a homomorphism into a small symmetric group.
\writeup (a) $56=2^3\cdot7$, $n_7\in\{1,8\}$. If $n_7=8$, the eight subgroups of order $7$ meet pairwise trivially and contain $48$ elements of order $7$. The remaining $8$ elements must contain a Sylow $2$-subgroup (which has $8$ elements, none of order $7$), so they form it: $n_2=1$, and that subgroup is normal. Either way there is a proper non-trivial normal subgroup.

(b) $200=2^3\cdot5^2$. $n_5$ divides $8$ and is $\equiv1\pmod5$; among $1,2,4,8$ only $1$ works. The Sylow $5$-subgroup is normal.

(c) $36=2^2\cdot3^2$, $n_3$ divides $4$ and is $\equiv1\pmod3$: $n_3\in\{1,4\}$. If $n_3=1$ we are done. If $n_3=4$, the conjugation action on the four Sylow $3$-subgroups is a homomorphism $\varphi:G\to\mathfrak S_4$. Its kernel is not $G$, since the action is transitive on $4$ points; and it is not $\{e\}$, since $|G|=36>24=|\mathfrak S_4|$. So $\ker\varphi$ is a proper non-trivial normal subgroup.
\end{solution}

\begin{solution}{ex:S5sub}
\idea A subgroup of small index $k$ yields a homomorphism to $\mathfrak S_k$ whose kernel is normal; but $\mathfrak S_5$ has very few normal subgroups.
\writeup Let $H\le\mathfrak S_5$ have index $k$. The action on the $k$ left cosets gives $\varphi:\mathfrak S_5\to\mathfrak S_k$; its kernel $K$ fixes the coset $H$, so $K\subseteq H$, and $|K|\ge120/k!$ because $\mathfrak S_5/K$ embeds in $\mathfrak S_k$. By Exercise~\ref{ex:normalSn}, $K\in\{\{e\},A_5,\mathfrak S_5\}$.

If $|H|=40$, $k=3$: $|K|\ge20$ and $|K|\le40$, but none of $1,60,120$ lies in $[20,40]$. If $|H|=30$, $k=4$: $|K|\ge5$ and $|K|\le30$; impossible again. If $|H|=15$: by Theorem~\ref{thm:pq} ($3\nmid4$), $H$ is cyclic, so $\mathfrak S_5$ would have an element of order $15$; but the largest order of an element of $\mathfrak S_5$ is $6$ (Exercise~\ref{ex:landau}: the seven partitions of $5$ give the orders $5,4,6,3,2,2,1$, never $15$).

(Subgroups of order $20$, $24$ and $60$ do exist: the normaliser of $\langle(1\,2\,3\,4\,5)\rangle$, the stabiliser of a point, and $A_5$.)
\end{solution}

\begin{solution}{ex:normalizer}
\idea Sylow's conjugacy theorem inside the smaller group $M$.
\writeup Let $g\in N_G(M)$. Then $gPg^{-1}\subseteq gMg^{-1}=M$, and $gPg^{-1}$ has order $p^a$, so it is a Sylow $p$-subgroup of $M$. But $P\trianglelefteq M$, so $P$ is the only Sylow $p$-subgroup of $M$ (Corollary~\ref{cor:sylownormal}, applied to $M$). Hence $gPg^{-1}=P$, i.e.\ $g\in M$. So $N_G(M)\subseteq M$, and the reverse inclusion always holds.
\end{solution}

\begin{solution}{ex:derived}
\idea The derived subgroup is the smallest normal subgroup with an abelian quotient, so the derived series descends as fast as any chain with abelian quotients can.
\writeup (a) Each $G^{(i+1)}$ is normal in $G^{(i)}$ with abelian quotient (Exercise~\ref{ex:commutator}), so if $G^{(k)}=\{e\}$, the derived series read backwards is a chain as in the definition. Conversely, let $(G_i)$ be such a chain. We show $G^{(i)}\subseteq G_{r-i}$ by induction: true for $i=0$; if $G^{(i)}\subseteq G_{r-i}$, then $G^{(i+1)}=(G^{(i)})'\subseteq(G_{r-i})'\subseteq G_{r-i-1}$, the last inclusion because $G_{r-i}/G_{r-i-1}$ is abelian (Exercise~\ref{ex:commutator}). So $G^{(r)}\subseteq G_0=\{e\}$.

(b) $\mathfrak S_n/A_n$ is abelian, so $\mathfrak S_n'\subseteq A_n$. Conversely, $[(a\,b),(a\,c)]=\bigl((a\,b)(a\,c)\bigr)^2=(a\,c\,b)^2=(a\,b\,c)$, so $\mathfrak S_n'$ contains all $3$-cycles, hence $A_n$ (Proposition~\ref{prop:gens}). So $\mathfrak S_4'=A_4$ and $\mathfrak S_5'=A_5$.

$A_4/V\cong\Z/3\Z$ is abelian, so $A_4'\subseteq V$. $A_4'$ is normal in $A_4$ (Exercise~\ref{ex:commutator}) and not trivial ($A_4$ is not abelian). A subgroup of order $2$ of $V$ is not normal in $A_4$: $(1\,2\,3)(1\,2)(3\,4)(1\,2\,3)^{-1}=(2\,3)(1\,4)$. So $A_4'=V$, and $V'=\{e\}$. The derived series of $\mathfrak S_4$ is $\mathfrak S_4\supset A_4\supset V\supset\{e\}$.

$A_5'$ is a non-trivial normal subgroup of the simple group $A_5$, so $A_5'=A_5$: the derived series of $\mathfrak S_5$ is $\mathfrak S_5\supset A_5=A_5=\cdots$, which never reaches $\{e\}$.
\end{solution}

\begin{solution}{ex:sylowGL2}
\idea The group acts on the $p+1$ lines of the plane $(\Z/p\Z)^2$; each Sylow subgroup is the set of ``shears'' fixing one line pointwise.
\writeup $|G|=(p^2-1)(p^2-p)=p(p-1)^2(p+1)$, whose $p$-part is $p$. The matrices $u_x=\begin{psmallmatrix}1&x\\0&1\end{psmallmatrix}$ satisfy $u_xu_y=u_{x+y}$, so $U\cong\Z/p\Z$ has order $p$: it is a Sylow $p$-subgroup.

For $g=\begin{psmallmatrix}a&b\\0&d\end{psmallmatrix}\in B$, a direct computation gives $gu_xg^{-1}=u_{ax/d}\in U$, so $B\subseteq N_G(U)$. Conversely, $u_xv=v$ with $x\ne0$ means $xv_2=0$, so the vectors fixed by all of $U$ form the line $L=\langle e_1\rangle$. If $gUg^{-1}=U$, then for every $u\in U$, $u\,(ge_1)=g\,(g^{-1}ug)\,e_1=ge_1$ since $g^{-1}ug\in U$; so $ge_1$ is fixed by $U$, $ge_1\in L$, and $g$ is upper triangular. Hence $N_G(U)=B$, of order $(p-1)^2p$, and $n_p=[G:B]=\frac{p(p-1)^2(p+1)}{p(p-1)^2}=p+1$.

Geometrically, $G$ acts on the $p+1$ lines through the origin of $(\Z/p\Z)^2$ (each of the $p^2-1$ non-zero vectors spans a line containing $p-1$ of them). $B$ is the stabiliser of the line $L$, and the conjugate $gUg^{-1}$ is the group of shears fixing the line $gL$ pointwise. So the Sylow $p$-subgroups correspond to the lines: $p+1$ of them, $\equiv1\pmod p$ as Sylow predicts.
\end{solution}

\begin{solution}{ex:order12}
\idea Split according to $n_3$. If $n_3=4$, the action on Sylow $3$-subgroups embeds $G$ in $\mathfrak S_4$. If $n_3=1$, $G$ is a semidirect product, and the possible actions are few.
\writeup $|G|=12=2^2\cdot3$, $n_3\in\{1,4\}$.

\emph{Case $n_3=4$.} The four subgroups of order $3$ have pairwise trivial intersections and contain $8$ elements of order $3$. Conjugation gives $\varphi:G\to\mathfrak S_4$. Its kernel is $\bigcap_iN_G(P_i)$, and $N_G(P_i)$ has index $n_3=4$, so order $3$: $N_G(P_i)=P_i$, and the intersection of the distinct $P_i$ is trivial. So $\varphi$ is injective and $G$ is isomorphic to a subgroup $H$ of order $12$ of $\mathfrak S_4$. As in the proof of Proposition~\ref{prop:a4}, $H$ has index $2$, so $g^2\in H$ for all $g$, and every $3$-cycle $g=(g^2)^2$ lies in $H$. The $3$-cycles generate $A_4$ (Proposition~\ref{prop:gens}), so $H=A_4$.

\emph{Case $n_3=1$.} Let $P=\langle a\rangle\trianglelefteq G$ have order $3$ and let $Q$ be a Sylow $2$-subgroup, of order $4$, so $Q\cong\Z/4\Z$ or $\Z/2\Z\times\Z/2\Z$ (Exercise~\ref{ex:order4}). As in Theorem~\ref{thm:pq}, $G\cong P\rtimes_\varphi Q$, where $\varphi:Q\to\operatorname{Aut}(P)=\{\mathrm{id},a\mapsto a^{-1}\}\cong\Z/2\Z$.
\begin{itemize}
\item $\varphi$ trivial: $G\cong\Z/3\Z\times Q$, that is $\Z/12\Z$ or $\Z/3\Z\times\Z/2\Z\times\Z/2\Z\cong\Z/2\Z\times\Z/6\Z$.
\item $Q=\langle b\rangle\cong\Z/4\Z$, $\varphi$ non-trivial: $\varphi_b$ must be inversion, so there is one such group, $\Z/3\Z\rtimes\Z/4\Z$.
\item $Q\cong\Z/2\Z\times\Z/2\Z$, $\varphi$ non-trivial: $\varphi$ is onto $\Z/2\Z$ with a kernel of order $2$, one of the three subgroups of order $2$ of $Q$. Every permutation of the three non-identity elements of $Q$ is an automorphism (the product of any two of them is the third), so for two such $\varphi,\varphi'$ there is $\alpha\in\operatorname{Aut}(Q)$ with $\varphi'=\varphi\circ\alpha$. Then $(n,h)\mapsto(n,\alpha^{-1}(h))$ is an isomorphism $P\rtimes_\varphi Q\to P\rtimes_{\varphi'}Q$: in $P\rtimes_{\varphi'}Q$,
\[
\begin{aligned}
{}&(n,\alpha^{-1}h)(n',\alpha^{-1}h')\\&\quad=\bigl(n\,\varphi'_{\alpha^{-1}h}(n'),\,\alpha^{-1}(hh')\bigr)\\
&\quad=\bigl(n\,\varphi_h(n'),\,\alpha^{-1}(hh')\bigr).
\end{aligned}
\]
 So there is one group. $D_6$ is such a group: with $r$ the rotation of order $6$ and $s$ a reflection, $\langle r^2\rangle$ is normal of order $3$, $Q=\{e,r^3,s,r^3s\}\cong(\Z/2\Z)^2$, and $s$ inverts $r^2$.
\end{itemize}

\emph{The five groups are different.} Only $\Z/12\Z$ has an element of order $12$; $\Z/2\Z\times\Z/6\Z$ is the other abelian one. Among the non-abelian ones, $A_4$ has $n_3=4$, while $D_6$ and $\Z/3\Z\rtimes\Z/4\Z$ have $n_3=1$; and the Sylow $2$-subgroups of $D_6$ are not cyclic, while $\Z/3\Z\rtimes\Z/4\Z$ has one isomorphic to $\Z/4\Z$ (and all are conjugate, hence isomorphic).
\end{solution}

\section*{Sources and further reading}

\begin{itemize}
\item L.~Sylow, \emph{Th\'eor\`emes sur les groupes de substitutions}, Math.\ Ann.\ 5 (1872), 584--594. The original theorems, stated for permutation groups.
\item H.~Wielandt, \emph{Ein Beweis f\"ur die Existenz der Sylowgruppen}, Arch.\ Math.\ 10 (1959), 401--402. The counting proof of Theorem~\ref{thm:sylow1}.
\item W.~Feit and J.~G.~Thompson, \emph{Solvability of groups of odd order}, Pacific J.\ Math.\ 13 (1963), 775--1029. The odd order theorem of the remark; formally verified in the Coq proof assistant in 2012.
\item H.~U.~Besche, B.~Eick and E.~A.~O'Brien, \emph{A millennium project: constructing small groups}, Internat.\ J.\ Algebra Comput.\ 12 (2002), 623--644. The groups of order up to $2000$, including the count for order $1024$.
\item S.~C.~Pohlig and M.~E.~Hellman, \emph{An improved algorithm for computing logarithms over $GF(p)$ and its cryptographic significance}, IEEE Trans.\ Inform.\ Theory 24 (1978), 106--110. The algorithm of the notebook.
\item D.~S.~Dummit and R.~M.~Foote, \emph{Abstract Algebra}, 3rd ed., Wiley, 2004, Chapters~4 and~6. Sylow theory with many worked classifications.
\end{itemize}
